% ============================================================
% SECȚIUNEA 2 — Suprafața API
% Fișiere sursă de referință:
%   backend/openapi-v1.json
%   backend/AIInterviewCoach.API/Controllers/
%   backend/AIInterviewCoach.API/RateLimiting/RateLimitingPolicies.cs
%   backend/AIInterviewCoach.API/Authorization/AuthorizationPolicies.cs
% ============================================================

\section{Suprafața API REST}
\label{sec:api_surface}

\subsection{Prezentare generală}

API-ul REST este expus prin ASP.NET Core Minimal Controllers și documentat
prin fișierul \texttt{backend/openapi-v1.json} (specificație OpenAPI 3.x).
Toate rutele au prefixul \texttt{/api}; fișierul de specificație a fost
generat cu \texttt{Swashbuckle.AspNetCore} și reflectă starea curentă a
codului sursă.

Aplicația definește \textbf{31 de endpoint-uri} organizate în cinci grupuri
(tag-uri OpenAPI): \texttt{Auth}, \texttt{Interviews}, \texttt{Problems},
\texttt{Submissions}, \texttt{AdminAuditLogs} și \texttt{AdminObservability}.
Endpoint-ul de generare hint-uri (\texttt{POST /api/problems/\{id\}/hints})
este prezent în controler dar nu apare în specificația OpenAPI curentă,
deoarece nu este adnotat cu \texttt{[ApiExplorerSettings]}.
% Sursă: Controllers/ProblemsController.cs:79

\subsection{Politici de autorizare și rate limiting}

Fiecare grup de endpoint-uri este protejat printr-o combinație de politici
de autorizare și de rate limiting.

\paragraph{Politici de autorizare}
% Sursă: API/Authorization/AuthorizationPolicies.cs
\begin{itemize}
  \item \texttt{InterviewerWorkspaceAccess} — rolul \texttt{Interviewer} sau \texttt{Admin}.
  \item \texttt{CandidateWorkspaceAccess} — rolul \texttt{Candidate} sau \texttt{Admin}.
  \item \texttt{AdminProblemManagement} — exclusiv \texttt{Admin}.
  \item \texttt{AdminAuditLogAccess} — exclusiv \texttt{Admin}.
  \item \texttt{AdminObservabilityAccess} — exclusiv \texttt{Admin}.
\end{itemize}

\paragraph{Politici de rate limiting}
% Sursă: API/RateLimiting/RateLimitingPolicies.cs
\begin{itemize}
  \item \texttt{AuthLogin} — fereastră fixă, 5 cereri/minut per adresă IP;
        protejează endpoint-urile de autentificare.
  \item \texttt{InterviewTokenAccess} — fereastră fixă, 20 cereri/minut per
        utilizator autentificat; protejează accesul prin token de interviu.
  \item \texttt{CandidateSubmissionFlow} — bucket de tokeni, 15 cereri/minut
        per utilizator autentificat; protejează submisiile și hint-urile.
  \item \texttt{AdminMutation} — fereastră fixă, 12 cereri/minut per utilizator
        autentificat; protejează operațiunile de scriere administrative.
\end{itemize}

\subsection{Tabel complet al endpoint-urilor}

Tabelul~\ref{tab:api_endpoints} listează toate cele 31 de endpoint-uri cu
metoda HTTP, calea, politica de autorizare aplicată și politica de rate
limiting (unde există).

% Legendă politici:
%   IW  = InterviewerWorkspaceAccess
%   CW  = CandidateWorkspaceAccess
%   APM = AdminProblemManagement
%   AAL = AdminAuditLogAccess
%   AOA = AdminObservabilityAccess
%   JWT = [Authorize] fără politică specifică (orice utilizator autentificat)
%   —   = fără autentificare obligatorie
%
%   AL  = AuthLogin
%   ITA = InterviewTokenAccess
%   CSF = CandidateSubmissionFlow
%   AM  = AdminMutation

\begin{table}[ht]
\centering
\caption{Endpoint-urile API REST}
\label{tab:api_endpoints}
\small
\begin{tabular}{llp{3.8cm}p{2.8cm}p{2.4cm}}
\hline
\textbf{Metodă} & \textbf{Grup} & \textbf{Cale} & \textbf{Autorizare} & \textbf{Rate limit} \\
\hline
\multicolumn{5}{l}{\textit{Auth — gestionarea sesiunii și a identității utilizatorului}} \\
\hline
POST & Auth & \texttt{/api/auth/register}       & —   & AuthLogin \\
POST & Auth & \texttt{/api/auth/login}           & —   & AuthLogin \\
GET  & Auth & \texttt{/api/auth/me}              & JWT & — \\
POST & Auth & \texttt{/api/auth/logout}          & —   & — \\
\hline
\multicolumn{5}{l}{\textit{Interviews — gestionarea interviurilor și a sesiunilor}} \\
\hline
POST & Interviews & \texttt{/api/interviews}                         & IW  & — \\
GET  & Interviews & \texttt{/api/interviews}                         & IW  & — \\
GET  & Interviews & \texttt{/api/interviews/\{id\}}                  & IW  & — \\
POST & Interviews & \texttt{/api/interviews/\{id\}/problems}         & IW  & — \\
GET  & Interviews & \texttt{/api/interviews/\{id\}/sessions}         & IW  & — \\
GET  & Interviews & \texttt{/api/interviews/sessions/\{sessionId\}}  & IW  & — \\
POST & Interviews & \texttt{/api/interviews/sessions/\{id\}/complete}& CW  & — \\
GET  & Interviews & \texttt{/api/interviews/token/\{token\}}         & —   & InterviewTokenAccess \\
POST & Interviews & \texttt{/api/interviews/token/\{token\}/start}   & CW  & InterviewTokenAccess \\
\hline
\multicolumn{5}{l}{\textit{Problems — gestionarea problemelor și a cazurilor de test}} \\
\hline
GET  & Problems & \texttt{/api/problems}                            & JWT & — \\
POST & Problems & \texttt{/api/problems}                            & APM & AdminMutation \\
GET  & Problems & \texttt{/api/problems/\{id\}}                     & JWT & — \\
PUT  & Problems & \texttt{/api/problems/\{id\}}                     & APM & AdminMutation \\
DELETE & Problems & \texttt{/api/problems/\{id\}}                   & APM & AdminMutation \\
GET  & Problems & \texttt{/api/problems/templates}                  & APM & — \\
POST & Problems & \texttt{/api/problems/signature-preview}          & APM & — \\
POST & Problems & \texttt{/api/problems/catalog/replace-with-starter-set} & APM & AdminMutation \\
POST & Problems & \texttt{/api/problems/\{problemId\}/testcases}    & APM & AdminMutation \\
GET  & Problems & \texttt{/api/problems/\{problemId\}/testcases}    & APM & — \\
POST & Problems & \texttt{/api/problems/\{id\}/hints}               & CW  & CandidateSubmissionFlow \\
\hline
\multicolumn{5}{l}{\textit{Submissions — trimiterea și consultarea soluțiilor}} \\
\hline
POST   & Submissions & \texttt{/api/submissions}                         & CW  & CandidateSubmissionFlow \\
GET    & Submissions & \texttt{/api/submissions/me}                      & CW  & CandidateSubmissionFlow \\
DELETE & Submissions & \texttt{/api/submissions/problem/\{problemId\}}   & CW  & CandidateSubmissionFlow \\
GET    & Submissions & \texttt{/api/submissions/session/\{sessionId\}}   & CW  & CandidateSubmissionFlow \\
DELETE & Submissions & \texttt{/api/submissions/session/\{sessionId\}}   & CW  & CandidateSubmissionFlow \\
\hline
\multicolumn{5}{l}{\textit{Admin — audit și observabilitate}} \\
\hline
GET & AdminAuditLogs      & \texttt{/api/admin/audit-logs}     & AAL & — \\
GET & AdminObservability  & \texttt{/api/admin/observability}  & AOA & — \\
\hline
\end{tabular}
\end{table}

\subsection{Convenții de răspuns}

\begin{itemize}
  \item \textbf{200 OK} — cerere procesată cu succes; corpul răspunsului
        conține DTO-ul corespunzător serializat JSON.
  \item \textbf{201 Created} — resursă nouă creată (interviuri, probleme).
  \item \textbf{204 No Content} — operațiune de ștergere sau deconectare
        finalizată fără corp de răspuns.
  \item \textbf{400 Bad Request} — validare eșuată
        (\texttt{ModelState.IsValid == false}) sau operațiune invalidă
        (\texttt{InvalidOperationException}).
  \item \textbf{401 Unauthorized} — token JWT absent sau expirat.
  \item \textbf{403 Forbidden} — utilizatorul este autentificat dar nu deține
        rolul necesar (\texttt{UnauthorizedAccessException}).
  \item \textbf{404 Not Found} — resursa solicitată nu există
        (\texttt{KeyNotFoundException}).
  \item \textbf{429 Too Many Requests} — depășirea limitei de cereri configurate.
  \item \textbf{500 Internal Server Error} — excepție neașteptată; mesajul
        de eroare este locat dar nu expus clientului.
\end{itemize}
% Sursă: API/Middleware/ExceptionHandlingMiddleware.cs

\subsection{Observații de design}

\paragraph{Endpoint-ul \texttt{GET /api/interviews/token/\{token\}}}
Singurul endpoint neautentificat din grupul Interviews. Permite candidaților
să vizualizeze detaliile unui interviu (titlu, probleme) înainte de autentificare,
folosind token-ul primit pe e-mail. Rate limiting-ul previne enumerarea
token-urilor prin forță brută.
% Sursă: Controllers/InterviewsController.cs:72

\paragraph{Endpoint-urile de ștergere a submisiilor}
Există două mecanisme de reset: \texttt{DELETE /api/submissions/problem/\{id\}}
șterge submisiile unui candidat pentru o singură problemă (în mod practice
sau în context de sesiune prin parametrul query \texttt{interviewSessionId}),
iar \texttt{DELETE /api/submissions/session/\{id\}} resetează întreaga sesiune
de interviu simultan.
% Sursă: Controllers/SubmissionsController.cs

\paragraph{Separarea rolurilor în \texttt{GET /api/problems}}
Endpoint-ul listează probleme diferit în funcție de rolul apelantului:
\texttt{Interviewer} vede propriile probleme, \texttt{Candidate} vede
problemele publice, \texttt{Admin} vede tot catalogul.
% Sursă: Application/Services/ProblemService.cs
